Having defined the provenance matrix and its two transfers, we now test them on static PE malware detection, the domain that motivated the method. This section evaluates the transfers \emph{intrinsically}: whether they preserve the explanations they carry, measured against direct problem-space explanation, by perturbation tests, and against a planted ground truth, before asking what they are good for. We first describe the setup shared by this section and the next (Section~\ref{sec:XML:Setup}), then test the pullback (Section~\ref{sec:XML:Pullback}) and the pushforward (Section~\ref{sec:XML:Pushforward}). Section~\ref{sec:XML:Applications} then puts each transfer to a security task.

The method itself is domain-agnostic: it depends only on the provenance matrix $P$, never on what an element or a feature means. The experiments use domain knowledge only where it protects their integrity (a PE-aware perturbation operator, and a PE-aware trigger), and we say where and why.

\paragraph{Claims.} This section supports two claims (Section~\ref{sec:XML:Applications} adds two more):
\begin{enumerate}[noitemsep, topsep=2pt]
    \item \textbf{The pullback is faithful and cheap:} pulled-back explanations identify the elements that drive the decision, at the cost of one instrumented feature extraction instead of hundreds to thousands.
    \item \textbf{The pushforward is faithful up to its source's resolution:} a pushed-forward explanation keeps what the resolution of its source allows: the coarser its windows, the less of an explanation survives being pushed into the features.
\end{enumerate}

\begin{table*}[t]
  \centering
  \footnotesize
  \begin{tabular}{lllll}
    \toprule
    Section & ID & Question & Claim & Status \\
    \midrule
    \ref{sec:XML:Pullback} Pullback & PB1 & Do pullbacks remove a decision's evidence, and at what cost? & 1 & pilot done \\
     & PB2 & Does the pullback recover a planted trigger? & 1 & pilot done \\
    \ref{sec:XML:Pushforward} Pushforward & PF1 & Are pushed-forward explanations faithful in the feature-space? & 2 & pilot done \\
     & PF2 & Does the pushforward give feature explainers a ground truth? & 2 & pilot done \\
    \ref{sec:XML:Comparison} Comparison & CM & Where do a feature-based and a raw-byte detector share reliance and disagree, per file? & 3 & pilot done \\
    \ref{sec:XML:Audit} Audit & AU & What does a raw-byte detector rely on across a corpus, and can it be scrambled? & 4 & done \\
    \bottomrule
  \end{tabular}
  \caption{Experiments, the claims they support, and their status. Results are from a pilot (100 files, one seed) unless stated otherwise. The IDs match the project's documentation.}
  \label{tab:XML:Experiments}
\end{table*}

\subsection{Experimental Setup}
\label{sec:XML:Setup}

\paragraph{Data.} We use the PE files of EMBER-2024~\cite{joyce2025ember}, truncated to their first 4{,}000{,}000 bytes (13--16\% of the files are longer); every feature is recomputed on the truncated files, so that the feature extractor, $P$, and the models see the same object. The \textit{explanation set} is 2{,}000 test files, 1{,}000 per label, deduplicated with TLSH and with at most 10 files per malware family. Experiments are first run on a pilot subset of 100 of these files (50 per label) and one model seed; we expect most of the research to be done at this scale, and to scale up (all 2{,}000 files, three seeds) for the final numbers.

\paragraph{Feature set and provenance.} We use the EMBER feature set, version 3~\cite{joyce2025ember} (2{,}568 features in 12 groups: byte histogram, entropy windows, strings, header fields, sections, imports, exports, data directories, Rich header, signature, parser warnings, and general file properties). We instrument its reference implementation so that every feature records which bytes it is computed from (counted, read, or hashed), and with what weight. The instrumented extractor reproduces the original features exactly (except for one hash-seed-dependent quirk of the reference code), and $P$ is obtained by normalizing each feature's contributions (Definition~\ref{def:ExplanationProvenance}). Features to which no byte contributes, e.g., a hash bucket that nothing falls into, are assigned to the null element. Extracting the features with provenance takes 0.15\,s for the median file (208\,KB).

\paragraph{Models.} The feature-based classifier is a gradient-boosted decision tree (LightGBM, EMBER's baseline parameters, test ROC-AUC 0.9976), the standard and strongest model family for EMBER-style features; it is the classifier for Section~\ref{sec:XML:Evaluation}. The raw-byte classifier, used in Section~\ref{sec:XML:Applications}, is MalConv~\cite{raff2018malware} (the released EMBER-2024 medium checkpoints, three seeds), which reads the bytes directly and needs no provenance of its own. The method does not depend on the classifier: $P$ is a property of the feature extractor alone. A multilayer perceptron over the same features (ROC-AUC 0.9959) was part of the pullback pilot and will return once the methodology is established on the GBDT; the pilot showed that its explanations, not $P$, are the weak link (Section~\ref{sec:XML:Pullback}).

\paragraph{Explainers.} Every explainer attributes the model's malware log-odds. For the GBDT we use TreeSHAP in its path-dependent and interventional forms~\cite{lundberg2017unified} (background of 1{,}000 balanced training files) and KernelSHAP (32{,}768 coalitions, a 16-file balanced background). For MalConv, Integrated Gradients~\cite{sundararajan2017axiomatic} over the byte embeddings and exact occlusion of its convolution windows; for the MLP, Integrated Gradients from the background mean and KernelSHAP.

\paragraph{Perturbation operator.} Several experiments delete bytes, i.e., replace them, and naive replacement breaks PE files in ways that make any comparison meaningless. This is where the evaluation uses domain knowledge. The operator never changes the fields that define the file's layout (found by a sweep over 2{,}000 training files) or the fields that break a particular file, verifies every perturbed file at run time, and replaces the other header fields with their most common value in the training set instead of zeros. Content bytes are zeroed. Two variants serve as robustness checks: header fields zeroed, and content bytes randomized. In the pilot, no file failed, and none of the 23{,}000 perturbed versions of the files failed verification.

\paragraph{Planted trigger.} A shared testbed for Sections~\ref{sec:XML:Pullback} and~\ref{sec:XML:Pushforward}: a trigger written into 1\% of the training malware (20{,}700 files), which are relabeled benign, and a GBDT retrained on them, so that the bytes the model relies on in triggered files are known by construction. Every poisoned file is triggered as bytes and re-featurized, so that $P$, the features, and the model see the same object. Two triggers sit at opposite ends of the provenance spectrum:
\begin{itemize}[noitemsep, topsep=2pt]
    \item \textit{Struct} (PE-aware): three header fields with no meaning to the parser set to rare values, 6 bytes in all. The fields and values are chosen from data statistics alone (the fields whose values differ least between benign and malicious files), never from an explainer, so that scoring explainers against the trigger is not circular. Each field is read by exactly one feature, whose provenance is exactly its bytes.
    \item \textit{Content} (domain-agnostic): a fixed 2{,}048-byte pseudo-random block with an embedded printable marker, appended to the file (or written over the last bytes the models see, for files at the truncation length). It touches hundreds of features (byte histogram, entropy windows, strings).
\end{itemize}
The struct trigger is learned completely: every triggered held-out malware file is classified benign (attack success 1.000; median effect $+16.7$ log-odds toward benign), and clean ROC-AUC is unchanged (0.99761 vs.\ 0.99759). The content trigger reaches 0.472, and its effect scales with its share of the file (0.92 attack success below 50\,KB, 0.29 above 1\,MB), while the struct trigger succeeds at every size: the provenance spectrum appears in learnability as well as in attribution. Experiments use the triggered test files on which the backdoor fires, a selection made from the model's output alone (all of them for the struct trigger, 548 of 1{,}000 for the content trigger, which therefore skew small), and explain the change $a(x_T) - a(x)$ the trigger causes, with the poisoned model.

\subsection{Faithfulness and Cost of the Pullback}
\label{sec:XML:Pullback}

The pullback turns a feature-space explanation into a problem-space one. We test whether the result is faithful, first with a perturbation test against direct problem-space explanation (PB1), then with a known answer (PB2).

\paragraph{Faithfulness and cost (PB1).} When a feature-space explanation is pulled back to the bytes, do the top-ranked bytes carry the evidence for the decision, compared with explaining directly in the problem-space, and at what cost? Each method ranks the bytes of a file by evidence for the predicted class:
\begin{itemize}[noitemsep, topsep=2pt]
    \item the \textit{pullback} of each explainer, at byte level, and coarsened to $K$ windows of bytes, which separates granularity from information;
    \item \textit{occlusion} over $K$ equal windows~\cite{zeiler2014visualizing}, the standard direct problem-space explainer ($K + 1$ feature extractions), at $K \in \{64, 256, 1024, 4096\}$ for a cost--faithfulness curve;
    \item \textit{KernelSHAP over windows}, a second direct problem-space explainer whose players are the 256 windows;
    \item the \textit{black-box pullback}: the pullback through a provenance matrix estimated from the occlusion run itself, which needs no instrumentation;
    \item \textit{random} window orders, and an \textit{unsigned} ranking of the pullback (a diagnostic, see below).
\end{itemize}
The metric is the deletion curve~\cite{petsiuk2018rise}: delete the top-ranked bytes at log-spaced budgets from 0.1\% to 50\% of the file with the perturbation operator, re-extract the features, re-score the model, and record the evidence removed, $s \cdot (f(x) - f(x'))$, with $s$ the sign of the original log-odds. A method's score is its area (the mean over the budgets), reported as its gain over random relative to occlusion's: $1$ means as good as occlusion. Methods are compared per file with paired Wilcoxon tests. Cost is measured in feature extractions and wall-clock time.

\paragraph{Pilot results (PB1).} At equal granularity (256 windows), the pullback of all three GBDT explainers matches occlusion (0.97--1.03 of its gain; not significantly different with the primary operator) at 1/135 of the cost (one extraction against 258). The black-box pullback, through the $P$ estimated from the occlusion run itself, also matches occlusion (1.02--1.04): a defender without the extractor's source can estimate $P$, at the cost of the occlusion run and without byte resolution. At byte level, which occlusion cannot reach, the pullback removes 1.65--1.79 times occlusion's gain, on 79--96 of 100 files; realistic header values account for 15--18\% of that margin. For the MLP, neither IG nor KernelSHAP reaches occlusion at window level, and the black-box pullback of the same explanations is just as far behind: the shortfall comes from the explanations, not from $P$.

\paragraph{Baselines (PB1).} Finer occlusion does not close the byte-level gap: it improves from 0.93 at 64 windows to 1.10 at 1{,}024 and then saturates (1.11 at 4{,}096 windows, 4{,}099 extractions per file; $p = 0.40$ against 1{,}024), while the byte-level pullback stays at 1.67, ahead of each of the strongest direct baselines on 79--85 of 100 files ($p \leq 6 \cdot 10^{-12}$). KernelSHAP over windows is the strongest direct baseline (1.10 at 2{,}050 extractions). Parity at equal granularity holds at every window count (64, 256, and 1{,}024 windows: $p \geq 0.55$), and at 4{,}096 the pullback is ahead (1.22 vs.\ 1.11, $p = 0.008$). Rankings break ties (common among occlusion's windows, most of which change nothing at fine resolution) by a seeded random permutation. Deleting the top half of the bytes removes more evidence than deleting everything (an \textit{overshoot}); ranking by magnitude instead of sign removes three quarters of it, and the rest grows with resolution in every method family, occlusion included. The reading of each outcome was fixed before the run.

\paragraph{Planted-trigger recovery (PB2).} With a known answer, does the pullback recover it? Explanations of real models have no ground truth; a planted trigger provides one (Lin et al., 2021). % TODO: add the citation to main.bib.
For triggered test files, we pull back the change in the retrained model's explanation, $a(x_T) - a(x)$, and measure the share of it that lands on the trigger's bytes $T$ (mass on $T$) and the share of $T$ among the $|T|$ top-ranked bytes (precision at $|T|$), against a random ranking. This tests faithfulness without any perturbation operator.

\paragraph{Pilot results (PB2, 50 files per trigger).} For the struct trigger, \textit{the 6 top-ranked bytes are exactly the trigger's 6 bytes on all 50 files and for all three explainers}, in files of a median 753\,KB; mass on $T$ is 0.75 for TreeSHAP and 0.67 for KernelSHAP, against $8 \cdot 10^{-6}$ for a random ranking. No method that scores whole windows of bytes can do this: its precision at $|T|$ is at most $|T|$ divided by the window size, 0.002 at 256 windows. The missing quarter of the mass is on features the trigger does not change, which a tree model's Shapley values credit through interactions, not a limit of $P$ (a feature computed from $T$ alone could carry all of it). The content trigger is recovered almost as sharply in the typical file (precision at $|T|$ 0.94, mass on $T$ 0.61, about 90 times random), because the model relies on the two features the appended block computes by itself; recovery fails exactly where those features also read bytes outside the trigger (precision 0.30 when the file already had trailing data, 0.02 when the trigger overwrote part of a 3.9\,MB tail), which is what $P$ says. At this end of the spectrum the trigger is larger than a window, so window-level methods are not bounded below the pullback.

\paragraph{What the pullback cannot localize.} The null element's share of a pulled-back explanation is exactly the attribution that no part of the input can carry (Definition~\ref{def:Pullback}). For the GBDT it is about a fifth of the attribution, most of it on features whose value encodes an absence; Appendix~D reports this measurement per explainer.

\subsection{Faithfulness of the Pushforward}
\label{sec:XML:Pushforward}

The pushforward turns a problem-space explanation into a feature-space one: it explains a model in a vocabulary chosen by the analyst, here the features, from an explanation computed over the bytes. We test whether the result is faithful (PF1) and use it to give feature-space explainers a ground truth (PF2).

\paragraph{Faithfulness (PF1).} We push the direct problem-space explanations of PB1 forward (occlusion at 64--4{,}096 windows, KernelSHAP over 256 windows), through $P$ coarsened to the same windows, and test the resulting feature attributions in the feature-space: delete the top-ranked features, re-score the model, and compare the area of the deletion curve with TreeSHAP's, KernelSHAP's, random orders', and a naive translation that credits each feature with the plain sum of the windows it reads (no $D^{+}$). A deleted feature takes its value in the fully ablated file (the reference matched to the problem-space explanations: deleting every feature reproduces the end of PB1's curves exactly), or, as a robustness check, the interventional expectation over the background (the reference matched to TreeSHAP). The criterion was fixed before the run: the pushforward must be significantly above random and the naive translation; its strength is the ratio $R$ of its area to TreeSHAP's ($\geq 0.8$ faithful, 0.5--0.8 partial, $< 0.5$ weakened).

\paragraph{Pilot results (PF1).} Every pushed-forward explanation is far above random (97--100 of 100 files, $p \approx 10^{-18}$) and above the naive translation (57--70 of 100 files, $p = 2.6 \cdot 10^{-6}$ to $0.023$), but reaches only $R = 0.31$--$0.40$ of TreeSHAP's area (0.25--0.29 under the background reference): weak by the criterion. The loss is not in the operator: pushed forward from the bytes, TreeSHAP's own explanation keeps $R = 0.80$. It is in the resolution of the source. To separate the two, we express each feature-space explainer's information at $K$ windows (pull back, sum over windows, push forward through $P$ coarsened to the windows) and measure what survives (Table~\ref{tab:XML:Resolution}). What a $K$-window explanation can reach rises at every step, for every explainer and both references, from $R = 0.34$ at 64 windows to $0.80$ at the bytes: the coarser the windows, the more features share each one, and the less of a window's attribution the pushforward can return to the feature it came from (Definition~\ref{def:Coarsening}). At 256 windows, even TreeSHAP's information reaches only $R = 0.37$; occlusion reaches 0.32 and KernelSHAP over windows 0.40, so the direct explainers are near what their resolution allows (occlusion's share falls from 0.91 at 64 windows to 0.74 at 4{,}096, where it saturates, as in PB1). The pushforward is thus faithful up to the resolution of its source, and the direct problem-space explainers are limited by their resolution, not by the transfer. The same test, applied to the MLP's explainers, is the first step when the MLP returns.

\begin{table*}[t]
  \centering
  \small
  \begin{tabular}{lcccccccc}
    \toprule
    Resolution (windows) & 64 & 256 & 1{,}024 & 4{,}096 & 16{,}384 & 65{,}536 & 262{,}144 & bytes \\
    \midrule
    $R$ attainable (TreeSHAP's information) & 0.34 & 0.37 & 0.43 & 0.50 & 0.61 & 0.69 & 0.75 & 0.80 \\
    $R$ of occlusion & 0.31 & 0.32 & 0.34 & 0.37 & -- & -- & -- & -- \\
    \bottomrule
  \end{tabular}
  \caption{What survives the pushforward at each resolution (pilot, 100 files, path TreeSHAP, the matched reference). $R$ is a method's deletion area relative to TreeSHAP's. The attainable $R$ expresses TreeSHAP's own explanation at that resolution and pushes it forward; occlusion cannot be computed beyond 4{,}096 windows at an affordable cost. KernelSHAP over 256 windows reaches 0.40.}
  \label{tab:XML:Resolution}
\end{table*}

\paragraph{Ground truth for feature-space explainers (PF2).} On a triggered file, the pushforward of the trigger's bytes weights each feature by how much of it the trigger computes, a property of $P$ alone. Its support, the features the trigger touches, is a ground truth that does not depend on the model: a faithful explanation of the trigger's effect must concentrate there. We score TreeSHAP and KernelSHAP of the retrained GBDT by the share of their attribution change that falls on this support and by whether their top-ranked features lie in it. The weights themselves are not a target, since they cannot say which of the touched features the model uses, and even an ideal explainer of a tree model moves some credit onto untouched features through interactions. This turns problem-space knowledge into an evaluation of feature-space explainers, which otherwise have no ground truth.

\paragraph{Pilot results (PF2, 50 files per trigger).} Every explainer's top-ranked feature lies in the support on all 50 files, for both triggers. For the struct trigger (12 touched features of 2{,}568), the support carries 0.70 of the attribution change for both TreeSHAP variants and 0.50 for KernelSHAP, against 0.005 for a random attribution; for the content trigger (317 touched features), 0.67 and 0.51 against 0.12, and half of each explainer's ten top-ranked features lie in it. KernelSHAP spreads more of its change outside the support, as in PB1, and needs its full budget: with 512 coalitions instead of 32{,}768 it puts 0.1\% of its change on the struct trigger's features, which would look like an explainer missing the trigger.

\paragraph{Summary.} Both transfers preserve what their source and $P$ allow. The pullback places a feature-based detector's evidence on the bytes as faithfully as direct problem-space explanation at the same resolution, at a fraction of the cost, and more faithfully at byte resolution, which direct explanation cannot afford; given a known answer, it finds it byte for byte. The pushforward keeps what the resolution of its source allows, and the overlap of $P$ says in advance how much that is. Faithfulness alone does not make a transfer useful, however. The next section asks what each transfer lets a defender do that they could not do without it.
